% TOP_PROBLEM: {"problemId": "problem.arnold-givental-lagrangian-intersection-conjecture", "canonicalTitle": "Arnold-Givental Lagrangian intersection conjecture", "releaseRank": 395, "primaryDomain": "geometry_topology", "primaryDomainLabel": "Geometry & topology", "displayStatus": "open_with_solved_subcases", "releaseStatus": "open_with_solved_subcases"}
% !TeX program = lualatex
% Definition and sources only; this file is not an LLM solution attempt.
\documentclass[11pt]{article}
\usepackage[margin=25mm]{geometry}
\usepackage{fontspec}
\setmainfont{Latin Modern Roman}
\usepackage{amsmath,amssymb,mathtools}
\usepackage{unicode-math}
\setmathfont{Latin Modern Math}
\usepackage{xurl,hyperref}
\hypersetup{colorlinks=true,urlcolor=blue}
\setcounter{secnumdepth}{0}
\setlength{\parindent}{0pt}
\setlength{\parskip}{5pt}
\setlength{\emergencystretch}{3em}
\begin{document}
\section*{\texorpdfstring{Arnold-Givental Lagrangian intersection conjecture}{Arnold-Givental Lagrangian intersection conjecture}}\label{395}
\subsection{Definitions and mathematical statement}
A real symplectic manifold is a compact symplectic $(M,\omega)$ with an involution $\tau$ satisfying $\tau^*\omega=-\omega$. Its nonempty fixed locus $L=\operatorname{Fix}(\tau)$ is Lagrangian: it has half the dimension of $M$ and $\omega|_{TL}=0$. A Hamiltonian diffeomorphism is the time-one map of a time-dependent vector field $X_{H_t}$ defined by $\iota_{X_{H_t}}\omega=-{\,\mathrm{d}} H_t$. If $L$ and $\phi(L)$ are transverse, the conjecture asserts
\[
 \#(L\cap\phi(L))\ge\sum_i\dim_{{\mathbb{F}}_2}H_i(L;{\mathbb{F}}_2).
\]
Compactness and transversality make the intersection finite. The real-locus hypothesis is part of the statement; it is not a lower bound for arbitrary Lagrangians or for non-Hamiltonian symplectomorphisms.
\par\smallskip\textit{Formulation and concept references:} \href{https://arxiv.org/html/1708.09628v1}{[S1]}.

\subsection{Short English statement}
Must every transverse Hamiltonian image of a compact real Lagrangian meet it in at least as many points as its total mod-two Betti number?
\subsection{Sources}
\begin{itemize}
\item[S1]Urs Frauenfelder, Arnold-Givental conjecture and moment Floer homology, 2017.
\url{https://arxiv.org/html/1708.09628v1}.
\textit{Formulation source.}
\end{itemize}
\end{document}
